% !TEX program = xelatex
% 驾驶安全事故预测与司机安全评价方案 —— 算法与模型说明
% 由 src/delivery/report.py 对应的内容生成；页脚注释已去除，仅保留页码。
\documentclass[UTF8,10pt]{ctexart}

\usepackage[a4paper,left=1.65cm,right=1.65cm,top=1.6cm,bottom=1.7cm]{geometry}
\usepackage[table]{xcolor}
\usepackage{booktabs}
\usepackage{array}
\usepackage{tabularx}
\usepackage{amsmath}
\usepackage{fancyhdr}

% ---------- 配色（与原 PDF 一致） ----------
\definecolor{heading}{HTML}{16344B}   % 标题深蓝
\definecolor{thead}{HTML}{E6EFF4}     % 表头浅蓝底
\definecolor{lineblue}{HTML}{A6BDCD}  % 表头下边框
\definecolor{linerow}{HTML}{D5DFE5}   % 行分隔线

% ---------- 页眉页脚：仅保留页码，页脚注释已去除 ----------
\pagestyle{fancy}
\fancyhf{}
\fancyfoot[R]{\small\thepage}
\renewcommand{\headrulewidth}{0pt}
\renewcommand{\footrulewidth}{0pt}

% ---------- 段落样式 ----------
\setlength{\parindent}{0pt}
\setlength{\parskip}{7pt}

% ---------- 自定义命令 ----------
\newcolumntype{Y}{>{\raggedright\arraybackslash}X}
\newcommand{\hdr}[1]{\par\vspace{0.4em}{\large\heiti\color{heading}#1}\par\vspace{0.4em}}
\newcommand{\secpage}[1]{\clearpage\hdr{#1}}
\newcommand{\thead}[1]{\textbf{\color{heading}#1}}

\begin{document}

% ================= 封面：方案概述与样本定义 =================
\begin{center}
{\fontsize{23}{32}\selectfont\heiti\color{heading}驾驶安全}\\[0.4em]
{\fontsize{23}{32}\selectfont\heiti\color{heading}事故预测与司机安全评价方案}\\[0.9em]
{\large\heiti\color{heading} \;\textperiodcentered\; 算法与模型说明}
\end{center}

\vspace{1.2em}

\hdr{方案概述}

本方案面向物流车队的事故预防和司机安全管理，使用组委会提供的车辆画像、风险事件、行驶轨迹及IMU数据，为500辆车输出未来40天事故风险预测，并建立可解释、可持续更新的司机安全评分体系。

事故预测采用正式目标与辅助风险目标的浅层LightGBM模型，通过训练期暴露量收缩、三随机种子平均和等权排名融合学习风险排序，再进行概率校准与预测期限换算。安全评分以历史高危事件、疲劳分心、驾驶稳定性、道路规则、管理配合和IMU佐证六个维度为基础，结合时间衰减、暴露校正及固定参考分布形成0至100分。

\begin{center}\small
\arrayrulecolor{lineblue}
\begin{tabularx}{\textwidth}{p{3.2cm} Y}
\rowcolor{thead}\thead{任务} & \thead{输出与用途} \\
\hline
事故预测 & 每辆车输出二分类\texttt{accident}及连续概率\texttt{risk\_prob}；按风险排序安排优先复核 \\
\hline
安全评价 & 每名司机输出安全分、排名、等级、分项贡献、前三项原因和管理建议 \\
\hline
落地与复现 & 提供数据处理、特征计算、训练预测、动态评分源码与安装运行说明 \\
\hline
\end{tabularx}
\end{center}

\hdr{样本与目标定义}

数据实际覆盖2026年6月1日至7月31日，共61天；\texttt{gpsno}作为车辆与司机的一一对应标识。正式正类定义为目标窗口内事故事件11803至少1条，或未遂事件11804至少2条。计数单位为事件记录，单条未遂可作为历史行为证据，但不单独构成正式正类。

最终预测截止日为2026年7月31日，预测区间为2026年8月1日至9月9日，共40天，计算采用左闭右开区间[2026-08-01, 2026-09-10)。只使用比赛提供的数据，不引入外部训练数据。

{总体流程：原始数据清洗 $\rightarrow$ 车辆日级汇总与时间标签 $\rightarrow$ 固定窗口特征 $\rightarrow$ 事故预测 / 六维安全评分 $\rightarrow$ 逐车结果与管理建议。历史评估用于检验风险区分能力；40天绝对概率的适用条件在预测方法与验证章节说明。}

% ================= 01 =================
\secpage{01\;|\;数据处理与特征工程}

\begin{center}\small
\arrayrulecolor{lineblue}
\begin{tabularx}{\textwidth}{p{2.4cm} p{2.6cm} Y}
\rowcolor{thead}\thead{数据来源} & \thead{规模} & \thead{处理结果} \\
\hline
车辆画像 & 500辆 & 字符串ID、唯一键和车辆集合校验 \\
\hline
风险事件 & 3,152,654条 & 逐事件表及完整车辆日期日历 \\
\hline
行驶轨迹 & 584,310,388条 & 里程、工时、速度、覆盖与质量日表 \\
\hline
车载IMU & 738,244,484条 & 各轴、模长、异常比例及缺失日表 \\
\hline
\end{tabularx}
\end{center}

清洗规则：使用DuckDB流式读取gzip并输出压缩Parquet。时间限制为[2026-06-01, 2026-08-01)，事件类型按编码识别。有效速度为0至160km/h；异常速度、坐标保留质量标记。距离由厘米转公里，负距离和负时长不进入暴露汇总。风险记录不会因位置异常而丢失标签，IMU缺失保持未知，不补造采样记录。

候选特征覆盖事件次数、暴露归一风险率、事件日密度、单日聚集性、近期趋势、历史事故与未遂、轨迹速度分布、IMU统计及传感器覆盖率。共构造795个候选特征，仅在最早训练窗口删除全空、常量和高度相关列（绝对相关系数$\geq$0.995），按特征族轮转保留100个核心列。

窗口一致性：核心列中93列使用固定14天观察窗，7列使用最近7天统计；最终预测与历史训练使用同一窗口含义。17个时间窗口共8,500行特征样本，用于不同期限的构造与审计；预测模型只使用14天观察、14天目标的样本。重复窗口来自同一批500辆车，不视为独立车辆。

辅助模型使用100个核心列和7个收缩风险率，覆盖历史高危、疲劳、分心、稳定性、道路、管理及盲区。对计数$n$、暴露$e$，定义 $r=(n+\tau\mu)/(e+\tau)$。$\mu$为训练期事件总量/暴露总量，$\tau$为训练期正暴露中位数。疲劳与分心按工时，其余按公里；零或未知暴露保留缺失。收缩用于降低低暴露下比率估计的波动。

防止信息泄漏：各窗口只使用截止日以前的信息，训练标签结束日不晚于预测截止日；收缩先验、裁剪阈值、填补值均只在训练期拟合。ID、日期、窗口元数据和未来目标不入模。半年画像统计截至7月31日，因此不作为早期回测的预测特征。

% ================= 02 =================
\secpage{02\;|\;事故预测模型与40天概率}

双目标建模：正式模型预测未来14天是否满足“事故$\geq$1或未遂$\geq$2”；辅助模型预测未来14天是否出现任一次事故或未遂。辅助目标提供更稠密的监督信号，两者使用同一时间边界，并统一按正式目标评价。模型采用浅层树与正则约束，以适应车辆数少、正式正类稀疏的特点。

\begin{center}\small
\arrayrulecolor{lineblue}
\begin{tabularx}{\textwidth}{p{4.4cm} Y Y}
\rowcolor{thead}\thead{参数} & \thead{正式目标模型} & \thead{辅助目标模型} \\
\hline
特征数 & 100 & 107 \\
\hline
树数 / 最大深度 / 叶子上限 & 63 / 2 / 7 & 100 / 3 / 7 \\
\hline
学习率 / 最小叶样本 & 0.03 / 40 & 0.03 / 30 \\
\hline
L1 / L2正则 & 0.5 / 5 & 0.5 / 10 \\
\hline
列采样 / 行采样 & 0.7 / 0.8 & 0.8 / 0.8 \\
\hline
类别权重 & 类别频数反比 & 正类乘以负正样本比平方根 \\
\hline
随机种子 & 2026 & 2026、42、137 \\
\hline
\end{tabularx}
\end{center}

正式模型先按训练期0.5\%与99.5\%分位裁剪，再用训练中位数填补并附加缺失指示；辅助模型使用LightGBM原生缺失分支。辅助训练按车辆出现窗口数的倒数归一基础权重，再叠加类别权重，减少重复窗口数量不同造成的偏置。

融合：三个辅助模型的概率先求平均；正式与辅助输出各自转换为完整车队的百分位$q_1$、$q_2$，风险分 $R=0.5q_1+0.5q_2$，同分采用平均名次。整批500辆车一起计算；推理检查唯一ID、完整车辆集合与统一截止日。最终训练使用截至7月31日已获得完整标签的14天样本。

概率校准：在6月30日截止的开发窗口拟合$C=1$的正则化Platt映射，得到14天概率$p_{14}$。二分类阈值在开发数据上满足召回率至少70\%的候选中最大化“召回率-假阳性率”，随后随预测期限作同一单调换算。

期限换算：假定事故和未遂是独立、平稳的Poisson计数，日强度分别为$\lambda_a$、$\lambda_n$，则 $P_t=1-\exp[-t(\lambda_a+\lambda_n)]\times(1+t\lambda_n)$。用$p_{14}$及历史事故/未遂条数占比求解强度，再取$t=40$得到\texttt{risk\_prob}；\texttt{accident}为\texttt{risk\_prob}达到换算阈值时的0/1标记。该外推保持排序，但40天绝对概率仍需未来真值验证。

% ================= 03 =================
\secpage{03\;|\;时间验证与预测效果}

\begin{center}\small
\arrayrulecolor{lineblue}
\begin{tabularx}{\textwidth}{p{2.8cm} p{3.4cm} Y}
\rowcolor{thead}\thead{环节} & \thead{特征截止日} & \thead{标签期与用途} \\
\hline
最早训练窗口 & 6月14日 & 6月15日至28日；特征筛选及初始监督样本 \\
\hline
开发评估 & 6月28日 / 6月30日 & 后续14天；参数确定与校准，训练标签须在截止日前可得 \\
\hline
历史回顾 & 7月17日 & 7月18日至31日；500辆车、28个正式正类 \\
\hline
\end{tabularx}
\end{center}

按时间先后划分样本，禁止随机拆分窗口行。每次评价只用当时可获得的训练标签。两个开发窗口标签重叠且各仅6个正类，不能视为独立重复实验。以下晚期结果属于已查看历史窗口的回顾性评估，不作为新的独立测试结论。

\begin{center}\small
\arrayrulecolor{lineblue}
\begin{tabularx}{\textwidth}{p{3.6cm} p{2.2cm} Y}
\rowcolor{thead}\thead{评价指标} & \thead{结果} & \thead{含义} \\
\hline
事故预测ROC-AUC & 0.8200 & 衡量正负车辆总体风险排序；比赛主要评价指标 \\
\hline
事故预测AP & 0.1577 & Average Precision，衡量稀疏正类的检索表现 \\
\hline
安全评分风险AUC & 0.6557 & 使用100减安全分作为风险方向 \\
\hline
\end{tabularx}
\end{center}

\begin{center}\small
\arrayrulecolor{lineblue}
\begin{tabularx}{\textwidth}{p{4.4cm} p{2.4cm} p{2.4cm}}
\rowcolor{thead}\thead{按预测风险优先复核} & \thead{捕获正类} & \thead{召回率} \\
\hline
前25辆 & 1 & 3.6\% \\
\hline
前50辆 & 9 & 32.1\% \\
\hline
前100辆 & 17 & 60.7\% \\
\hline
\end{tabularx}
\end{center}

高风险最前段仍存在不稳定性，前25辆仅捕获1个正类。管理使用时应结合评分明细与事件证据确定复核对象，不能把较高总体AUC理解为最前端车辆一定发生事故。

{敏感性检验：随机20\%车辆IMU输入缺失时，与完整输入预测的排序Spearman为0.9533，前50辆重合率为86\%。}

{敏感性检验：全部车辆IMU输入缺失时，与完整输入预测的排序Spearman为0.9828，前50辆重合率为82\%。}

{三个辅助随机种子间排序Spearman约0.975至0.979。上述检验描述输入或随机种子扰动下的稳定性，不代表未来性能保证。开发正类少、同车窗口重叠和传感器覆盖随时间变化均限制结论外推；当前没有8月以后真值，真实40天风险区分与概率校准尚未验证。}

% ================= 04 =================
\secpage{04\;|\;安全评分公式与维度权重}

安全评分采用截至评分日的固定14天观察窗，分数越高表示观测行为风险越低。评分与事故概率独立呈现：前者解释近期行为及管理重点，后者估计目标期事件风险，两者不能直接互换。

第一步，时间衰减。事件距截止日$a$天时权重为$2^{-a/14}$，维度事件强度 $N=\sum[\text{事件严重度}\times2^{-a/14}]$。里程、工时和活跃日按相同14天半衰期衰减，得到暴露$E$。事件按发生时刻计算，日级暴露按当日中点近似。

第二步，暴露收缩。五个行为维度计算 $r=(N+\tau\mu)/(E+\tau)$。参考$\mu$为训练期总事件强度/总暴露，$\tau$为训练期正暴露中位数。历史、稳定性、道路按每千公里，疲劳分心按每百小时，管理按实际移动车辆日。暴露不足时向参考总体收缩，避免将少行驶或少上报直接视作安全。

第三步，参考映射。用6月1日至14日的固定参考分布，将风险率映射为惩罚指数$P$：$P=\text{上限}\times[\text{参考值小于}r\text{的数量}+0.5\times\text{等于}r\text{的数量}]/\text{参考样本数}$。五个行为维度上限100，IMU上限60。固定参考保证不同日期的分数可比较。

\begin{center}\small
\arrayrulecolor{lineblue}
\begin{tabularx}{\textwidth}{p{4.2cm} p{2.4cm} Y}
\rowcolor{thead}\thead{维度} & \thead{权重} & \thead{暴露口径} \\
\hline
历史事故与未遂 & 26.60\% & 每千公里 \\
\hline
疲劳与分心 & 13.97\% & 每百小时 \\
\hline
驾驶稳定性 & 26.60\% & 每千公里 \\
\hline
道路规则 & 7.38\% & 每千公里 \\
\hline
管理配合 & 15.45\% & 每活跃日 \\
\hline
IMU佐证 & 10.00\% & 采样记录及日级统计 \\
\hline
\end{tabularx}
\end{center}

第四步，加权合成。维度贡献$D=w\times P$，安全分$S=100-\sum D$，范围为0至100。贡献包含经验先验收缩，解释中同时披露实际记录及收缩成分。IMU缺失时行为贡献为0，另列$U=30\times w_{IMU}$及保守下界$\max(0,S-U)$，表示证据不充分，不把缺失直接记为违规。

{权重以正式与辅助目标的成对逻辑排序损失等权训练，加入向六维业务先验（35\%、20\%、20\%、10\%、10\%、5\%）的L2收缩（系数2），由SLSQP求解。约束权重非负且和为1、与先验单维偏差不超过0.10、历史权重最高、IMU不超过10\%、疲劳及稳定性不低于道路规则。训练使用6月14日评分和后续14天标签，参考分布及权重在动态应用中固定。}

% ================= 05 =================
\secpage{05\;|\;行为细则、IMU佐证与可信度}

事件严重度表示相对危险程度，并非直接扣分。每条事件只归属一个行为维度，经过时间衰减、暴露校正和参考映射后再计算贡献；24类事件均纳入规则。

\begin{center}\small
\arrayrulecolor{lineblue}
\begin{tabularx}{\textwidth}{p{3.2cm} Y}
\rowcolor{thead}\thead{维度} & \thead{事件类型与严重度} \\
\hline
历史事故与未遂 & 事故10；未遂事故7 \\
\hline
疲劳与分心 & 疲劳（闭眼）5；困倦4.5；打哈欠3.5；看手机3.5；注意力分散3；打电话3；频繁低头3；抽烟2 \\
\hline
驾驶稳定性 & 前碰撞预警5；车距过近3；长时间压线2.5；左/右车道偏移各1.5；后/右盲区报警各2 \\
\hline
道路规则 & 路口超速2.5；主路弯道超速3；危险路段超速3；匝道超速3；高速长下坡超速4 \\
\hline
管理配合 & 摄像头遮挡2；摄像头角度扭转1.5 \\
\hline
\end{tabularx}
\end{center}

IMU佐证：不假定设备坐标轴对应车体方向，使用加速度模长和角速度模长。令$f_a$为加速度模长不在0.5至1.5范围的记录比例，$f_g$为角速度模长大于20的记录比例，$d_v$为EMS与GPS速度差的非空日均值。速度差可得时 $r_{IMU}=0.60f_a+0.30f_g+0.10\times\min(d_v/20,1)$；不可得时去掉该项并将其余项除以0.90。

IMU指标是风险佐证，不作为事故标签。异常采样点数可能来自同一次持续动作，解释中不将其称为独立危险事件次数。没有IMU记录时，该维度为未知，并通过可信度与不确定性下界披露。

数据可信度$Q$由覆盖率、有效率和暴露充分度组成，范围0至100：$Q=25\times\text{轨迹日覆盖率}+15\times\text{IMU日覆盖率}+15\times\text{速度有效率}+15\times\text{坐标有效率}+10\times\min(\text{衰减里程}/1000,1)+10\times\min(\text{衰减工时}/100,1)+10$。日覆盖率按14天计算并截断到[0,1]；最后10分表示事件数据源可读取，不代表所有风险事件均被捕获。

$Q\geq80$为高可信度，$60\leq Q<80$为中可信度，$Q<60$为低可信度。低可信度优先排查设备与数据覆盖，不直接据分数奖惩。每名司机展示贡献最大的前三个维度，包括记录数量、惩罚指数、扣分贡献及对应建议，明细贡献之和应等于100减安全分。

{权重训练只有4个正式正类，学习到的权重仍受小样本影响。模型通过业务先验、非负约束、IMU扣分上限和人工复核流程约束应用范围。}

% ================= 06 =================
\secpage{06\;|\;风险分级与动态运营}

按安全分降序形成排名，名次越小越安全；同分按\texttt{gpsno}确定稳定顺序。A/B/C/D采用20\%/40\%/30\%/10\%的固定容量，即500辆车分别为100/200/150/50辆。分级服务于有限管理资源的分配，不代表绝对事故概率；车队整体改善应观察固定参考安全分及行为率，而非等级人数。

\begin{center}\small
\arrayrulecolor{lineblue}
\begin{tabularx}{\textwidth}{p{2.2cm} p{2.2cm} p{2.4cm} Y}
\rowcolor{thead}\thead{等级} & \thead{车辆数} & \thead{后续正类} & \thead{观察率 / Wilson 95\%区间} \\
\hline
A & 100 & 1 & 1.0\% / 0.2\%--5.4\% \\
\hline
B & 200 & 12 & 6.0\% / 3.5\%--10.2\% \\
\hline
C & 150 & 9 & 6.0\% / 3.2\%--11.0\% \\
\hline
D & 50 & 6 & 12.0\% / 5.6\%--23.8\% \\
\hline
\end{tabularx}
\end{center}

以上为7月17日评分对应7月18日至31日目标的回顾结果。各等级观察率为1\%、6\%、6\%、12\%，但置信区间明显重叠，尚不足以证明等级间差异稳定。Wilson区间只反映当前样本的二项抽样不确定性，不能覆盖时间漂移与模型选择偏差。

\begin{center}\small
\arrayrulecolor{lineblue}
\begin{tabularx}{\textwidth}{p{3.2cm} Y}
\rowcolor{thead}\thead{等级 / 状态} & \thead{行动与周期} \\
\hline
A：常规观察 & 维持监测；复核数据与事件后实施正向激励 \\
\hline
B：行为复盘 & 每月复盘前三项贡献及可改善行为 \\
\hline
C：定向培训 & 针对主要维度安排培训，每周检查事件与风险率 \\
\hline
D：优先干预 & 优先复核具体事件，形成培训、整改、复查闭环 \\
\hline
低可信度 & 设备、安装角度和上报覆盖检查优先；暂不自动奖惩 \\
\hline
\end{tabularx}
\end{center}

更新机制：每天、每周或每月触发相同评分流程，始终取截至当日的14天滚动观察窗，使用固定先验、参考分布及权重。新行为进入窗口，旧行为衰减或移出，分数随之变化；执行频率不改变观察窗口长度。

提交结果截至7月31日，并提供相对7月24日的7日分数变化；7月24日至31日连续8天评分用于验证动态计算。管理者同时关注安全分、风险率、扣分原因与可信度，司机查看自身分数、排名和改进重点。评分配置所需标签截至7月14日，动态应用不得早于该日期；超过已有数据覆盖日需先接入新数据。

% ================= 07 =================
\secpage{07\;|\;运行复现与提交文件}

运行环境：Python 3.13，CPU执行，无需GPU。主要依赖DuckDB、pandas、NumPy、LightGBM、scikit-learn和SciPy；PDF生成使用ReportLab。代码包提供锁定依赖、安装运行说明以及数据预处理、特征、训练、预测、评分和校验源码。

{\small 输入目录：}
\begin{verbatim}
data/raw/vehicle_profile.xlsx
data/raw/risk_events.tsv.gz
data/raw/trajectory/*.gz（30个分片）
data/raw/imu/*.gz（30个分片）
\end{verbatim}

{\small 安装与运行：}
\begin{verbatim}
python -m venv .venv
.venv/Scripts/python.exe -m pip install -r requirements-lock.txt
.venv/Scripts/python.exe -m src.reproduce_submission --clean --rebuild-features --package
\end{verbatim}

{\small 已有清洗表和特征：}
\begin{verbatim}
.venv/Scripts/python.exe -m src.reproduce_submission --package
\end{verbatim}

{\small 按日期更新评分：}
\begin{verbatim}
.venv/Scripts/python.exe -m src.scoring.enhanced --as-of 2026-07-31
\end{verbatim}

{\small 提交文件校验：}
\begin{verbatim}
.venv/Scripts/python.exe -m src.delivery.validate submission/final
\end{verbatim}

Linux/macOS将解释器路径替换为.venv/bin/python。预测读取固定模型配置，重新拟合各训练期先验、校准器和最终模型；评分读取固定权重及参考分布。全量数据处理采用流式读取，运行时间取决于硬件与磁盘。

\begin{center}\small
\arrayrulecolor{lineblue}
\begin{tabularx}{\textwidth}{p{3.8cm} Y}
\rowcolor{thead}\thead{提交文件} & \thead{内容及格式} \\
\hline
\texttt{forecast\_result.csv} & 500行；\texttt{gpsno}唯一；\texttt{accident}为0/1整数；\texttt{risk\_prob}为有限[0,1]概率 \\
\hline
\texttt{safety\_scores.csv} & 同一500名司机；安全分[0,100]、排名、等级、六维贡献、可信度、前三项原因、7日变化及建议 \\
\hline
\texttt{method\_report.pdf} & 数据处理、特征、预测、评分、验证与运营方案说明 \\
\hline
\texttt{source\_code.zip} & 源代码、模型配置、依赖与安装运行说明；比赛数据由授权使用者另行提供 \\
\hline
\end{tabularx}
\end{center}

结果CSV使用UTF-8无BOM与LF换行。验收覆盖车辆集合、预测取值、评分贡献对账、排名与等级一致性。源码包从基础清洗表独立重建特征、训练及评分，已验证两份结果CSV逐字节一致；该复现范围不包含再次扫描全部原始传感器文件。

\end{document}
